% GENERATED FILE. Edit canonical lessons, then run npm run generate:books.
% canonical-source-hash: fnv1a64:1de0531ad491741e
% canonical-lessons: TE-C128-gaddam, TE-C128-manikattu, TE-C128-moceyi, TE-C128-carmam, TE-C128-raktam

\chapter{Chin, Wrist, Elbow, Skin, Blood}
\label{ch:te-chin-wrist-elbow-skin-blood}
\hlchaptermodality{128}

\begin{chapteropening}
\textbf{By the end of this chapter:} \emph{I can say the chin, a wrist, an elbow, the skin and blood.}

Complete the last lesson of chapter 128: i can say the chin, a wrist, an elbow, the skin and blood.
\end{chapteropening}

\section[\texorpdfstring{\te{గడ్డం} (gaḍḍaṁ)}{గడ్డం (gaḍḍaṁ)}]{\te{గడ్డం} (gaḍḍaṁ) — the chin; a beard}

\label{lesson:TE-C128-gaddam}



\begin{quote}
Before the new one: say the Telugu for an eyebrow, then the Telugu for a cheek.
\end{quote}

\subsection*{You'll want to know: \te{గడ్డం}}
\textbf{\te{గడ్డం}} — \emph{gaḍḍaṁ} — \textquotedblleft{}the chin; a beard\textquotedblright{}.

\begin{grammarlens}[title={What you've built}]
One more word for this chapter.
\end{grammarlens}

\subsection*{Your turn}
\begin{itemize}
\raggedright
  \item \emph{Say it:} \emph{gaḍḍaṁ}
  \item \emph{Say it:} \emph{gaḍḍaṁ}, once more
  \item \emph{Say it:} say \emph{gaḍḍaṁ}
  \item \emph{Recall:} say the Telugu for an eyebrow, then the Telugu for a cheek, then say \emph{gaḍḍaṁ} again
\end{itemize}

\begin{tcolorbox}[breakable,colback=teal!4,colframe=teal!35!black,title={Before you move on}]
What does \te{గడ్డం} mean? (\textquotedblleft{}The chin; a beard\textquotedblright{}.) Say it once more.
\end{tcolorbox}
\section[\texorpdfstring{\te{మణికట్టు} (maṇikaṭṭu)}{మణికట్టు (maṇikaṭṭu)}]{\te{మణికట్టు} (maṇikaṭṭu) — a wrist}

\label{lesson:TE-C128-manikattu}



\begin{quote}
Before the new one: say the Telugu for a cheek, then the Telugu for the chin; a beard.
\end{quote}

\subsection*{You'll want to know: \te{మణికట్టు}}
\textbf{\te{మణికట్టు}} — \emph{maṇikaṭṭu} — \textquotedblleft{}a wrist\textquotedblright{}.

\begin{grammarlens}[title={What you've built}]
One more word for this chapter.
\end{grammarlens}

\subsection*{Your turn}
\begin{itemize}
\raggedright
  \item \emph{Say it:} \emph{maṇikaṭṭu}
  \item \emph{Say it:} \emph{maṇikaṭṭu}, once more
  \item \emph{Say it:} say \emph{maṇikaṭṭu}
  \item \emph{Recall:} say the Telugu for a cheek, then the Telugu for the chin; a beard, then say \emph{maṇikaṭṭu} again
\end{itemize}

\begin{tcolorbox}[breakable,colback=teal!4,colframe=teal!35!black,title={Before you move on}]
What does \te{మణికట్టు} mean? (\textquotedblleft{}A wrist\textquotedblright{}.) Say it once more.
\end{tcolorbox}
\section[\texorpdfstring{\te{మోచేయి} (mōcēyi)}{మోచేయి (mōcēyi)}]{\te{మోచేయి} (mōcēyi) — an elbow}

\label{lesson:TE-C128-moceyi}



\begin{quote}
Before the new one: say the Telugu for the chin; a beard, then the Telugu for a wrist.
\end{quote}

\subsection*{You'll want to know: \te{మోచేయి}}
\textbf{\te{మోచేయి}} — \emph{mōcēyi} — \textquotedblleft{}an elbow\textquotedblright{}.

\begin{grammarlens}[title={What you've built}]
One more word for this chapter.
\end{grammarlens}

\subsection*{Your turn}
\begin{itemize}
\raggedright
  \item \emph{Say it:} \emph{mōcēyi}
  \item \emph{Say it:} \emph{mōcēyi}, once more
  \item \emph{Say it:} say \emph{mōcēyi}
  \item \emph{Recall:} say the Telugu for the chin; a beard, then the Telugu for a wrist, then say \emph{mōcēyi} again
\end{itemize}

\begin{tcolorbox}[breakable,colback=teal!4,colframe=teal!35!black,title={Before you move on}]
What does \te{మోచేయి} mean? (\textquotedblleft{}An elbow\textquotedblright{}.) Say it once more.
\end{tcolorbox}
\section[\texorpdfstring{\te{చర్మం} (carmaṁ)}{చర్మం (carmaṁ)}]{\te{చర్మం} (carmaṁ) — the skin}

\label{lesson:TE-C128-carmam}



\begin{quote}
Before the new one: say the Telugu for a wrist, then the Telugu for an elbow.
\end{quote}

\subsection*{You'll want to know: \te{చర్మం}}
\textbf{\te{చర్మం}} — \emph{carmaṁ} — \textquotedblleft{}the skin\textquotedblright{}.

\begin{grammarlens}[title={What you've built}]
One more word for this chapter.
\end{grammarlens}

\subsection*{Your turn}
\begin{itemize}
\raggedright
  \item \emph{Say it:} \emph{carmaṁ}
  \item \emph{Say it:} \emph{carmaṁ}, once more
  \item \emph{Say it:} say \emph{carmaṁ}
  \item \emph{Recall:} say the Telugu for a wrist, then the Telugu for an elbow, then say \emph{carmaṁ} again
\end{itemize}

\begin{tcolorbox}[breakable,colback=teal!4,colframe=teal!35!black,title={Before you move on}]
What does \te{చర్మం} mean? (\textquotedblleft{}The skin\textquotedblright{}.) Say it once more.
\end{tcolorbox}
\section[\texorpdfstring{\te{రక్తం} (raktaṁ)}{రక్తం (raktaṁ)}]{\te{రక్తం} (raktaṁ) — blood}

\label{lesson:TE-C128-raktam}



\begin{quote}
Before the new one: say the Telugu for an elbow, then the Telugu for the skin.
\end{quote}

\subsection*{You'll want to know: \te{రక్తం}}
\textbf{\te{రక్తం}} — \emph{raktaṁ} — \textquotedblleft{}blood\textquotedblright{}.

\begin{grammarlens}[title={What you've built}]
One more word for this chapter.
\end{grammarlens}

\subsection*{Your turn}
\begin{itemize}
\raggedright
  \item \emph{Say it:} \emph{raktaṁ}
  \item \emph{Say it:} \emph{raktaṁ}, once more
  \item \emph{Say it:} say \emph{raktaṁ}
  \item \emph{Recall:} say the Telugu for an elbow, then the Telugu for the skin, then say \emph{raktaṁ} again
\end{itemize}

\begin{tcolorbox}[breakable,colback=teal!4,colframe=teal!35!black,title={Before you move on}]
What does \te{రక్తం} mean? (\textquotedblleft{}Blood\textquotedblright{}.) Say it once more.
\end{tcolorbox}
